\subsectionoptional{长度与夹角的度量}
我们已经把第一节中关于解集的结论翻译成了几何语言，以便更好地理解这些集合。
但我们必须当心，不要被自己的术语所误导\Dash 给 \( \Re^k \) 中形如
\( \set{\vec{p}+t\vec{v}\suchthat t\in\Re} \) 与
\( \set{\vec{p}+t\vec{v}+s\vec{w}\suchthat t,s\in\Re} \)
的子集标注为"直线"与"平面"，并不能使它们表现得像我们过去经验中的
直线和平面那样。
确切地说，我们必须确保这些名称与这些集合相称。
虽然我们无法证明这些集合满足我们的直觉\Dash 关于直觉我们无法证明任何东西\Dash
但在本小节中我们将观察到，
一个在 \( \Re^2 \) 与 \( \Re^3 \) 中熟悉的结果，当推广到任意的 \( \Re^n \) 时，
支持了直线是直的、平面是平的这一观念。
具体而言，我们将给出两个 $\Re^n$ 向量在它们所生成的平面内的夹角的定义，
来看看如何在一个"平面"内做欧几里得几何。

\begin{definition} \label{df:Length}
%<*df:Length>
向量 \( \vec{v}\in\Re^n \) 的\definend{长度}\index{vector!length}\index{length of a vector}
% (or \definend{norm}\index{vector!norm}\index{norm})
是其各分量的平方和的平方根。
\begin{equation*}
  \absval{\vec{v}\,}=\sqrt{v_1^2+\cdots+v_n^2}
\end{equation*}
%</df:Length>
\end{definition}

\begin{remark}
这是勾股定理的一个自然推广。
一篇经典的引导性讨论见 \cite{MathematicsPlausReason}。
\end{remark}

对任意非零的 $\vec{v}$，向量 $\vec{v}/\absval{\vec{v}}$ 的长度为一。
我们说这一做法将
$\vec{v}$ \definend{规范化}\index{normalize, vector}\index{vector!normalize}为长度一。

我们可以用它来得到两个向量之间夹角的公式。
考虑 \( \Re^3 \) 中两个向量，它们互不为对方的倍数
\begin{center}
  \includegraphics{gr/mp/ch1.21}
\end{center}
（互为倍数的特殊情形将在下文看到并非例外）。
它们确定一个二维平面\Dash
例如，把它们放到标准位置，取由原点与两个端点所形成的平面。
在该平面中考虑以 \( \vec{u} \)、\( \vec{v} \) 与 \( \vec{u}-\vec{v} \) 为边的三角形。
\begin{center}
  \includegraphics{gr/mp/ch1.22}
\end{center}
应用余弦定理：
$\absval{\vec{u}-\vec{v}\,}^2
  =
  \absval{\vec{u}\,}^2+\absval{\vec{v}\,}^2-
    2\,\absval{\vec{u}\,}\,\absval{\vec{v}\,}\cos\theta$
其中 \( \theta \) 是两个向量之间的夹角。
左边给出
\begin{multline*}
(u_1-v_1)^2+(u_2-v_2)^2+(u_3-v_3)^2  \\
     =(u_1^2-2u_1v_1+v_1^2)+(u_2^2-2u_2v_2+v_2^2)+(u_3^2-2u_3v_3+v_3^2)
\end{multline*}
而右边给出下式。
\begin{equation*}
(u_1^2+u_2^2+u_3^2)+(v_1^2+v_2^2+v_3^2)-
     2\,\absval{\vec{u}\,}\,\absval{\vec{v}\,}\cos\theta
\end{equation*}
消去平方项 $u_1^2$、\ldots、$v_3^2$ 并除以 $2$，便得到夹角的公式。
\begin{equation*}
  \theta
  =
  \arccos(\,\frac{u_1v_1+u_2v_2+u_3v_3}{
         \absval{\vec{u}\,}\,\absval{\vec{v}\,} }\,)
\end{equation*}

在更高维的空间中我们无法像上面那样作图，
但可以改用解析的方法来进行论证。
首先，分子的形式是清楚的；它来自 \( (u_i-v_i)^2 \) 的中间项。

\begin{definition} \label{df:DotProduct}
%<*df:DotProduct>
两个 \( n \) 分量实向量的\definend{点积}\index{vector!dot product}\index{dot product}
（或\definend{内积}\index{inner product}
或\definend{标量积}\index{scalar product}）
是其分量的线性组合。
\begin{equation*}
  \vec{u}\dotprod\vec{v}=\lincombo{u}{v}
\end{equation*}
%</df:DotProduct>
\end{definition}
\noindent 注意两个向量的点积是一个实数，
而不是向量，并且
只有当两个向量的分量个数相同时，点积才有定义。
还要注意点积与长度相关：
\( \vec{u}\dotprod\vec{u}=u_1u_1+\cdots+u_nu_n=\absval{\vec{u}\,}^2 \)。

\begin{remark}
有些作者要求第一个向量为行向量，
第二个向量为列向量。
我们将不这样严格，而是允许在两个列向量之间进行点积运算。
\end{remark}

我们仍在用解析的方法推理，但以图形为指导，
用下面的定理来论证：由 \( \Re^n \) 中构成
\( \vec{u} \)、\( \vec{v} \) 与 \( \vec{u}+\vec{v} \) 主体的线段所形成的三角形，
位于 \( \vec{u} \) 与
\( \vec{v} \) 生成的 \( \Re^n \) 的平面子集中（见下图）。

\begin{theorem}[三角不等式] \label{th:TriangleInequality}
\index{Triangle Inequality}
%<*th:TriangleInequality>
对任意 \( \vec{u},\vec{v}\in\Re^n \)，
\begin{equation*}
  \absval{\vec{u}+\vec{v}\,}\leq\absval{\vec{u}\,}+\absval{\vec{v}\,}
\end{equation*}
当且仅当其中一个向量是另一个向量的非负标量倍数时取等号。
%</th:TriangleInequality>
\end{theorem}

这就是人们熟悉的那句话的来源：
"两点之间直线最短。"
\begin{center}
  \includegraphics{gr/mp/ch1.23}
\end{center}

\begin{proof}
（我们将用到点积的一些尚未验证的代数性质，例如
\( \vec{u}\dotprod(\vec{a}+\vec{b})
   =\vec{u}\dotprod\vec{a}+\vec{u}\dotprod\vec{b} \)
以及 $\vec{u}\dotprod\vec{v}=\vec{v}\dotprod\vec{u}$。
参见 \nearbyexercise{exer:AlgPropsDotProd}。）
%<*pf:TriangleInequality0>
由于所有的数都是正的，
该不等式成立当且仅当两边平方后的不等式成立。
\begin{align*}
  \absval{\vec{u}+\vec{v}\,}^2
  &\leq(\,\absval{\vec{u}\,}+\absval{\vec{v}\,}\,)^2                            \\
  (\,\vec{u}+\vec{v}\,)\dotprod(\,\vec{u}+\vec{v}\,)
  &\leq\absval{\vec{u}\,}^2+2\,\absval{\vec{u}\,}\,\absval{\vec{v}\,}
     +\absval{\vec{v}\,}^2                                         \\
  \vec{u}\dotprod\vec{u}+\vec{u}\dotprod\vec{v}
     +\vec{v}\dotprod\vec{u}+\vec{v}\dotprod\vec{v}
  &\leq\vec{u}\dotprod\vec{u}+2\,\absval{\vec{u}\,}\,\absval{\vec{v}\,}
    +\vec{v}\dotprod\vec{v}                                          \\
  2\,\vec{u}\dotprod\vec{v}
  &\leq 2\,\absval{\vec{u}\,}\,\absval{\vec{v}\,}
\end{align*}
%</pf:TriangleInequality0>
%<*pf:TriangleInequality1>
而这一个不等式又当且仅当
将两边乘以非负数
\( \absval{\vec{u}\,} \)
与 \( \absval{\vec{v}\,} \)
所得到的
\begin{equation*}
  2\,(\,\absval{\vec{v}\,}\,\vec{u}\,)\dotprod(\,\absval{\vec{u}\,}\,\vec{v}\,)
  \leq
  2\,\absval{\vec{u}\,}^2\,\absval{\vec{v}\,}^2
\end{equation*}
再改写为
\begin{equation*}
  0
  \leq
  \absval{\vec{u}\,}^2\,\absval{\vec{v}\,}^2
   -2\,(\,\absval{\vec{v}\,}\,\vec{u}\,)\dotprod(\,\absval{\vec{u}\,}\,\vec{v}\,)
  +\absval{\vec{u}\,}^2\,\absval{\vec{v}\,}^2
\end{equation*}
这一关系成立。
但通过因式分解可知它是成立的
\begin{equation*}
  0\leq
  (\,\absval{\vec{u}\,}\,\vec{v}-\absval{\vec{v}\,}\,\vec{u}\,)\dotprod
  (\,\absval{\vec{u}\,}\,\vec{v}-\absval{\vec{v}\,}\,\vec{u}\,)
\end{equation*}
因为它只是说向量
\( \absval{\vec{u}\,}\,\vec{v}-\absval{\vec{v}\,}\,\vec{u}\, \)
的长度的平方不是负数。
%</pf:TriangleInequality1>
%<*pf:TriangleInequality2>
至于等号，它成立当且仅当
\( \absval{\vec{u}\,}\,\vec{v}-\absval{\vec{v}\,}\,\vec{u} \) 为 \( \zero \)。
容易验证：
\( \absval{\vec{u}\,}\,\vec{v}=\absval{\vec{v}\,}\,\vec{u}\, \)
当且仅当一个向量是另一个向量的非负实标量倍数。
%</pf:TriangleInequality2>
\end{proof}

这一结果支持了如下直觉：即使在更高维的空间中，直线也是直的，平面也是平的。
从定义出发我们容易验证，线性曲面具有这样的性质：
该曲面中的任意两点，
它们之间的线段仍包含在该曲面中。
但若线性曲面不是平的，那就可能出现捷径。
\begin{center}
  \includegraphics{gr/mp/ch1.24}
\end{center}
由于三角不等式表明，在任意 $\Re^n$ 中两个端点之间的最短路线
就是连接它们的线段，
所以线性曲面没有弯曲。

回到角度度量的定义。
三角不等式证明的核心是
\( \vec{u}\dotprod\vec{v}\leq \absval{\vec{u}\,}\,\absval{\vec{v}\,} \)
这一行。
我们或许会想知道，是否有些向量对以下述方式满足该不等式：~虽然 \( \vec{u}\dotprod\vec{v} \)
是一个较大的数，其绝对值大于右边，
但它是负的较大数。
下面的结论表明这种情况不会发生。

\begin{corollary}[Cauchy–Schwarz 不等式]
\label{th:CauchySchwarz}
\index{Cauchy-Schwarz Inequality}\index{Schwarz Inequality}
%<*co:CauchySchwarzInequality>
对任意 \( \vec{u},\vec{v}\in\Re^n \)，
\begin{equation*}
   \absval{\,\vec{u}\dotprod\vec{v}\,}
   \leq
   \absval{\,\vec{u}\,}\,\absval{\vec{v}\,}
\end{equation*}
当且仅当一个向量是另一个向量的标量倍数时取等号。
%</co:CauchySchwarzInequality>
\end{corollary}

\begin{proof}
%<*pf:CauchySchwarzInequality>
三角不等式的证明表明
\( \vec{u}\dotprod\vec{v}\leq \absval{\vec{u}\,}\,\absval{\vec{v}\,} \)，所以
若 $\vec{u}\dotprod\vec{v}$ 为正或为零，则证明已完成。
若 \( \vec{u}\dotprod\vec{v} \) 为负，则下式成立。
\begin{equation*}
 \absval{\,\vec{u}\dotprod\vec{v}\,}
   =-(\,\vec{u}\dotprod\vec{v}\,)
   =(-\vec{u}\,)\dotprod\vec{v}
   \leq
   \absval{-\vec{u}\,}\,\absval{\vec{v}\,}
   =\absval{\vec{u}\,}\,\absval{\vec{v}\,}
\end{equation*}
等号成立的条件见 \nearbyexercise{exer:EqOfCauchySchwarz}。
%</pf:CauchySchwarzInequality>
\end{proof}

Cauchy–Schwarz 不等式向我们保证下一个定义是有意义的，
因为该分式的绝对值小于或等于一。

\begin{definition}
%<*df:Angle>
两个非零向量 \( \vec{u},\vec{v}\in\Re^n \) 之间的\definend{夹角}\index{vector!angle}\index{angle}为
\begin{equation*}
  \theta
  =
  \arccos(\,\frac{\vec{u}\dotprod\vec{v}}{
                  \absval{\vec{u}\,}\,\absval{\vec{v}\,} }\,)
\end{equation*}
（若其中有一个是零向量，则我们取夹角为直角）。
%</df:Angle>
\end{definition}

\begin{corollary} \label{co:VectorsOrthogonalIffDoTProductZero}
%<*co:VectorsOrthogonalIffDoTProductZero>
\( \Re^n \) 中的向量
是正交的，\index{vector!orthogonal}\index{orthogonal}
即垂直的，\index{vector!perpendicular}\index{perpendicular}
当且仅当它们的点积为零。
它们平行当且仅当它们的点积等于
它们长度的乘积。
%</co:VectorsOrthogonalIffDoTProductZero>
\end{corollary}

\begin{example}
这些向量是正交的。
\begin{center}
  \begin{minipage}{0.6in} % to vertically center graphic
    \includegraphics{gr/mp/ch1.25}
  \end{minipage}
  \qquad
  $\colvec[r]{1 \\ -1}\dotprod\colvec[r]{1 \\ 1}=0$
\end{center}
我们把箭头画在离开标准位置的地方，
但尽管如此，这两个向量仍是正交的。
\end{example}

\begin{example}
本小节开头给出的 \( \Re^3 \) 夹角公式
是该定义的一个特例。
在这两个向量之间
\begin{center}
  \includegraphics{gr/mp/ch1.26}
\end{center}
夹角为
\begin{equation*}
   \arccos(\frac{(1)(0)+(1)(3)+(0)(2)}{\sqrt{1^2+1^2+0^2}\sqrt{0^2+3^2+2^2}})
   =\arccos(\frac{3}{\sqrt{2}\sqrt{13}})
\end{equation*}
约为 $0.94~\text{弧度}$。
注意这两个向量并不正交。
虽然 \( yz \)-平面看起来可能与
\( xy \)-平面垂直，但实际上这两个平面只是在如下的弱意义上如此：
每个平面中都存在与另一平面中所有向量正交的向量。
但每个平面中并非所有向量都与另一平面中的所有向量正交。
\end{example}




